\section{Limitations and Future Work}

The scenario design in this work is informed by the common distinction between authentication factors — something the user \textit{knows}, something the user \textit{has}, and something the user \textit{is} — as defined in NIST SP 800-63B~\cite{nist80063b}. However, the three scenarios do not constitute single-, two-, and three-factor authentication in the formal sense of this standard. NIST SP 800-63B defines possession-based factors as physical or cryptographic authenticators, such as hardware tokens or cryptographic keys bound to a device. A booking code does not meet this definition because it is a knowledge-based secret that can be copied, shared, or forwarded. S2 therefore does not introduce a second authentication factor; instead, its security gain comes from contextual enforcement at the moment of entry rather than from a distinct possession factor.

A second limitation concerns the scale of the evaluation. Each scenario was executed with 20 bookings and 20 access attempts under controlled and identical infrastructure conditions. While this allows for a direct comparison between scenarios, it does not reflect production-scale traffic patterns, concurrent access attempts, or variable load conditions. The observed cost and resource usage values should therefore be interpreted as indicative rather than representative of real-world deployments.

Finally, the evaluation focuses exclusively on the Access Service (R4) as the measurement point. Costs and resource usage associated with supporting services such as Amazon Rekognition, the Nuki Web API, or the database layer are included in the cost metric M1 but are not individually isolated. This means that the contribution of each external dependency to the overall cost difference between scenarios cannot be precisely attributed from the current measurements alone.